% This file is part of the modular source for Stratified Relational Workspace Theory.
% !TeX root = ../Stratified_Relational_Workspace_Theory.tex
\theorypart{XII}{Identity Through Time and Replacement}

Identity is treated as invariance under a structured correspondence, not as
literal persistence of tokens or carriers. Exact correspondences form a
groupoid. Approximate identity is measured by structural, operational,
causal, and reconstructive defects that must be reported separately before
they are combined.

\topic{38}{Temporal Realization Maps}

Let
\[
\mathfrak R=(G,\mathbf C,D,\rho,\mathcal I,\mathcal K)
\]
and
\[
\mathfrak R'=(G',\mathbf C',D',\rho',\mathcal I',\mathcal K')
\]
be two realizations. A correspondence begins with a graph isomorphism
\(\varphi:G\to G'\). After relabeling by \(\varphi\), let
\[
J_i:C_i\to C_i'
\]
be unitary local identifications and define
\[
J_V=\bigoplus_{i\in V}J_i,
\qquad
J_E=\bigoplus_{e\in E}J_{t(e)}.
\]

\begin{definition}[Exact structural correspondence]\label{def:exact-structural-correspondence}
The correspondence is structurally exact when
\[
\boxed{D'J_V=J_ED.}
\]
It is operation exact for \(a\in\mathcal A\) when
\[
\boxed{\rho_V'(a)J_V=J_V\rho_V(a).}
\]
\end{definition}

\begin{theorem}[Exact replacement isomorphism]\label{thm:exact-replacement-isomorphism}
Under an exact structural correspondence,
\[
\boxed{
\Gamma(D')=J_V\Gamma(D),
\qquad
L'J_V=J_VL.
}
\]
Consequently, every spectral workspace subspace is transported exactly:
\[
P_I(L')=J_VP_I(L)J_V^*
\]
for every spectral interval \(I\). If the correspondence is also operation
exact, then the transported workspace actions are intertwined exactly.
\end{theorem}

\begin{proof}
For every \(c\),
\[
D'J_Vc=J_EDc.
\]
Because \(J_E\) is invertible, \(Dc=0\) if and only if \(D'J_Vc=0\), proving
the kernel statement. Taking adjoints of the exact covariance identity gives
\[
D'^*J_E=J_VD^*.
\]
Therefore
\[
L'J_V=D'^*D'J_V=D'^*J_ED=J_VD^*D=J_VL.
\]
The spectral and operation claims follow by functional calculus and the
intertwining assumptions.
\end{proof}

Exact correspondences compose and invert. They therefore form a groupoid of
realizations. The abstract exact workspace is an isomorphism class in this
groupoid.

\topic{39}{Structural Identity Defect}

For an approximate correspondence define
\[
\boxed{
\Delta_D=D'J_V-J_ED,
\qquad
\delta_D=\|\Delta_D\|.
}
\]
If the graph weights are unchanged, define edgewise comparison defects
\[
\delta_e
=
\|H_e'J_{s(e)}-J_{t(e)}H_e\|.
\]
The discrepancy operator has one changed source block per edge, so
\[
\boxed{
\delta_D^2
\leq
\max_{i\in V}
\sum_{e:s(e)=i}w_e\delta_e^2
\leq
\sum_{e\in E}w_e\delta_e^2.
}
\]
When a single node is replaced, every nonincident edge has \(\delta_e=0\).
The replacement defect is therefore locally measurable even though its effect
on the global workspace is controlled by the global spectral gap.

The corresponding carried-Laplacian defect obeys
\[
\boxed{
\eta_L
:=
\|J_V^*L'J_V-L\|
\leq
2\|D\|\delta_D+\delta_D^2.
}
\]
This is Proposition~\ref{prop:temporal-laplacian-drift} with the time indices
removed.

Let \(P\) and \(P'\) be rank-\(d\) workspace projectors selected from \(L\)
and \(L'\), and let \(g>0\) be the separating source gap.

\begin{theorem}[Approximate structural replacement]\label{thm:approx-structural-replacement}
If
\[
\eta_L<\frac g2,
\]
then
\[
\boxed{
q_{\mathrm{str}}
:=
\|P'-J_VPJ_V^*\|
\leq
\frac{\eta_L}{g-\eta_L}
\leq
\frac{2\eta_L}{g}.
}
\]
If \(q_{\mathrm{str}}<1\), then
\[
B=P'J_VP:\ran P\longrightarrow\ran P'
\]
is bijective and satisfies
\[
\sigma_{\min}(B)
\geq
\sqrt{1-q_{\mathrm{str}}^2}.
\]
Its polar factor
\[
\boxed{
\beta
=B(B^*B)^{-1/2}
}
\]
is a canonical unitary workspace correspondence. Moreover,
\[
\boxed{
\|\beta-J_VP\|
\leq
q_{\mathrm{str}}
+
1-\sqrt{1-q_{\mathrm{str}}^2}.
}
\]
\end{theorem}

\begin{proof}
The projector estimate is Theorem~\ref{thm:spectral-stability} applied after
carrying \(L'\) back by \(J_V\). For unit \(w\in\ran P\),
\[
\|(I-P')J_Vw\|
\leq q_{\mathrm{str}}.
\]
Orthogonality gives
\[
\|Bw\|^2
=1-\|(I-P')J_Vw\|^2
\geq1-q_{\mathrm{str}}^2.
\]
Thus \(B\) is injective, hence bijective because the source and target have
the same finite dimension. Its singular values lie in
\([\sqrt{1-q_{\mathrm{str}}^2},1]\), so polar normalization moves \(B\) by at
most \(1-\sqrt{1-q_{\mathrm{str}}^2}\). Finally,
\[
\|B-J_VP\|=\|(I-P')J_VP\|\leq q_{\mathrm{str}}.
\]
\end{proof}

\subtopic{39.1 Transport of workspace operations}

Write
\[
A=\rho_V(a),
\qquad
A'=\rho_V'(a),
\]
and define
\[
\delta_\rho(a)=\|A'J_V-J_VA\|,
\qquad
\ell_a=\|(I-P)AP\|.
\]
Let
\[
A_{\mathcal W}=PAP,
\qquad
A_{\mathcal W'}'=P'A'P'.
\]

\begin{proposition}[Operation transport bound]\label{prop:operation-transport-bound}
For \(B=P'J_VP\),
\[
\boxed{
\|A_{\mathcal W'}'B-BA_{\mathcal W}\|
\leq
\|A'\|q_{\mathrm{str}}
+\delta_\rho(a)
+\ell_a.
}
\]
For the unitary polar correspondence \(\beta\),
\[
\boxed{
\begin{aligned}
\|A_{\mathcal W'}'\beta-\beta A_{\mathcal W}\|
\leq
&\|A'\|q_{\mathrm{str}}
+\delta_\rho(a)
+\ell_a\\
&+
\bigl(\|A_{\mathcal W'}'\|+\|A_{\mathcal W}\|\bigr)
\left(1-\sqrt{1-q_{\mathrm{str}}^2}\right).
\end{aligned}
}
\]
\end{proposition}

\begin{proof}
Insert and subtract \(P'A'J_VP\) and \(P'J_VAP\):
\[
\begin{aligned}
A_{\mathcal W'}'B-BA_{\mathcal W}
={}&P'A'(P'-I)J_VP\\
&+P'(A'J_V-J_VA)P\\
&+P'J_V(I-P)AP.
\end{aligned}
\]
Bound the three terms. Replacing \(B\) by \(\beta\) adds
\[
A_{\mathcal W'}'(\beta-B)-(\beta-B)A_{\mathcal W},
\]
which gives the second estimate.
\end{proof}

\topic{40}{Behavioral Identity Defect}

Structural correspondence does not determine operational identity. Let
\[
\{\iota_{w,b}:(w,b)\in\mathcal D\}
\quad\text{and}\quad
\{\iota'_{\beta w,\chi b}:(w,b)\in\mathcal D\}
\]
be matched feasible implementation-intervention families in the sense of
Part~VIII. Here \(b\) denotes declared background variables and \(\chi\) is
the background correspondence. Define
\[
K(w,b)
=
\mathbb P^{\iota_{w,b}}(Y^+\in\cdot),
\]
\[
K'(\beta w,\chi b)
=
\mathbb P^{\prime\,\iota'_{\beta w,\chi b}}(Y^{\prime+}\in\cdot).
\]
For a chosen divergence \(D_Y\), define
\[
\boxed{
\delta_{\mathrm{causal}}(\beta)
=
\sup_{(w,b)\in\mathcal D}
D_Y\left(
K(w,b),
K'(\beta w,\chi b)
\right).
}
\]
The observational response is included only when the null intervention belongs
to the declared family. This definition compares actual implementation
interventions rather than formally assigning an arbitrary projection
coordinate.

A reconstructive handoff may also be imperfect. Let
\[
\delta_{\mathrm{rec}}(\beta)
=
\mathbb E\bigl[d_{\mathcal W'}(W'_{\mathrm{init}},\beta W)\bigr]
\]
measure replacement initialization error. Corollary~\ref{cor:handoff-initialization}
provides one explicit upper bound on this term.

\subtopic{40.1 Structural identity does not imply causal identity}

Take two realizations with identical
\[
G,
\quad D,
\quad L,
\quad P,
\quad \rho.
\]
In the first, let \(Y=W\). In the second, let \(Y'=1-W\) for a binary
workspace coordinate. Then
\[
q_{\mathrm{str}}=0,
\qquad
\delta_\rho(a)=0
\]
for every operation, while the observational and interventional outcome
kernels differ maximally under the identity correspondence. Therefore no
purely structural theorem can prove causal identity without an additional
kernel-matching assumption.

\topic{41}{Identity Distance and Replacement Invariance}

For an admissible correspondence \(\mathfrak J\), define the operation defect
on a declared unit ball \(\mathcal A_1\) by
\[
\delta_{\mathrm{op}}(\mathfrak J)
=
\sup_{a\in\mathcal A_1}
\|A_{\mathcal W'}'\beta-\beta A_{\mathcal W}\|.
\]
The primary certification object is the identity profile
\[
\boxed{
\mathbf d_{\mathrm{id}}(\mathfrak J)
=
\left(
q_{\mathrm{str}},
\delta_{\mathrm{op}}(\mathfrak J),
\delta_{\mathrm{causal}}(\beta),
\delta_{\mathrm{rec}}(\beta);
\mu^{\mathrm{keep}}
\right).
}
\]
The first four coordinates must remain below declared tolerances, while the
non-collapse margin \(\mu^{\mathrm{keep}}\) must remain above a declared lower
bound. This vector, rather than a weighted average, determines whether a
specific identity claim passes.

A composite identity cost is
\[
\boxed{
\operatorname{Cost}(\mathfrak J)
=
a\,q_{\mathrm{str}}
+b\,\delta_{\mathrm{op}}(\mathfrak J)
+c\,\delta_{\mathrm{causal}}(\beta)
+d\,\delta_{\mathrm{rec}}(\beta),
}
\]
where \(a,b,c,d>0\) are declared scales. Define
\[
\boxed{
d_{\mathrm{id}}(\mathfrak R,\mathfrak R')
=
\inf_{\mathfrak J}
\operatorname{Cost}(\mathfrak J).
}
\]
This quantity is a task-relative optimization surrogate until the
admissible correspondence class, intervention matching, divergence, and scales
are fixed. It does not replace the coordinatewise identity profile or the
non-collapse gate.

\subtopic{41.1 Replacement invariance theorem \(T_{14}\) }

\begin{theorem}[Quantitative replacement invariance]\label{thm:replacement-invariance}
Let \(\mathfrak J\) be a unitary graph correspondence with diagram defect
\(\delta_D\). Define
\[
\eta_L=2\|D\|\delta_D+\delta_D^2
\]
and assume
\[
\eta_L<\frac g2.
\]
Set
\[
q_*
=
\frac{\eta_L}{g-\eta_L}.
\]
Then there exists a unitary workspace correspondence \(\beta\) for which
\[
q_{\mathrm{str}}\leq q_*.
\]
Moreover,
\[
\boxed{
\begin{aligned}
\delta_{\mathrm{op}}(\mathfrak J)
\leq
\sup_{a\in\mathcal A_1}
\Bigl[&
\|\rho_V'(a)\|q_*
+\delta_\rho(a)
+\ell_a\\
&+
\bigl(\|A_{\mathcal W'}'\|+\|A_{\mathcal W}\|\bigr)
\left(1-\sqrt{1-q_*^2}\right)
\Bigr].
\end{aligned}
}
\]
If, in addition,
\[
\delta_{\mathrm{causal}}(\beta)\leq\delta_c,
\qquad
\delta_{\mathrm{rec}}(\beta)\leq\delta_r,
\]
then
\[
\boxed{
\begin{aligned}
d_{\mathrm{id}}(\mathfrak R,\mathfrak R')
\leq
&a\,q_*+b\,\delta_{\mathrm{op}}(\mathfrak J)\\
&+c\,\delta_c+d\,\delta_r.
\end{aligned}
}
\]
In particular, if all local structural, operation, closure, causal, and
reconstruction defects are \(O(\delta_{\mathrm{rep}})\), the relevant
operator norms are uniformly bounded, and \(g\geq g_0>0\), then
\[
\boxed{
d_{\mathrm{id}}(\mathfrak R,\mathfrak R')
\leq
C_1\frac{\delta_{\mathrm{rep}}}{g_0}
+C_2\delta_{\mathrm{rep}}
+O\!\left(\frac{\delta_{\mathrm{rep}}^2}{g_0}\right).
}
\]
After a fixed normalization of the non-spectral terms, this recovers the
schematic target
\[
d_{\mathrm{id}}
\lesssim
C\frac{\delta_{\mathrm{rep}}}{g_0}.
\]
\end{theorem}

\begin{proof}
The structural bound and existence of \(\beta\) follow from
Theorem~\ref{thm:approx-structural-replacement}. The operation estimate is
Proposition~\ref{prop:operation-transport-bound} with
\(q_{\mathrm{str}}\leq q_*\). Insert these estimates and the assumed causal
and reconstructive bounds into the definition of
\(\operatorname{Cost}(\mathfrak J)\). The asymptotic statement follows from
\[
q_*=O\!\left(\frac{\delta_D}{g_0}
+\frac{\delta_D^2}{g_0}\right)
\]
and
\[
1-\sqrt{1-q_*^2}=O(q_*^2).
\]
\end{proof}

\begin{corollary}[Preservation of coalition frame bounds]
\label{cor:replacement-frame-bound}
Identify the old and new abstract cores with \(\ran P\) and \(\ran P'\)
through their chosen core parameterizations. Let \(S\) be a declared coalition,
let
\[
A_S:\ran P\to\mathsf Y_S,
\qquad
A_S':\ran P'\to\mathsf Y_S'
\]
be the resulting accessible coalition encoders, and let \(K_S\) be a unitary
identification of their observation spaces. Suppose
\[
\|A_S'\beta-K_SA_S\|
\leq
\delta_{S}^{\mathrm{frame}}
<
\sqrt{\alpha_S}.
\]
Then the replacement coalition remains reconstructive and its new lower frame
bound satisfies
\[
\boxed{
\alpha_S'
\geq
\left(
\sqrt{\alpha_S}-\delta_{S}^{\mathrm{frame}}
\right)^2.
}
\]
Consequently, \(k\)-failure resilience survives whenever the bound holds
uniformly for every coalition remaining after at most \(k\) erasures.
\end{corollary}

\begin{proof}
For \(w'\in\ran P'\), write \(w'=\beta w\). Since \(\beta\) and \(K_S\) are
unitary,
\[
\begin{aligned}
\|A_S'w'\|
&\geq
\|K_SA_Sw\|
-
\|(A_S'\beta-K_SA_S)w\|\\
&\geq
\left(
\sqrt{\alpha_S}-\delta_{S}^{\mathrm{frame}}
\right)\|w'\|.
\end{aligned}
\]
Square the inequality.
\end{proof}

\begin{corollary}[Exact carrier replacement]\label{cor:exact-carrier-replacement}
If
\[
\delta_D=0,
\qquad
\delta_\rho(a)=0,
\qquad
\ell_a=0,
\qquad
\delta_c=0,
\qquad
\delta_r=0,
\]
then
\[
d_{\mathrm{id}}(\mathfrak R,\mathfrak R')=0.
\]
Thus a participant, model, artifact, or notation system may be replaced without
changing the abstract workspace when the replacement realizes the same
structural, operational, causal, and reconstructive organization.
\end{corollary}

\subtopic{41.2 Temporal chains and identity paths}

For a sequence of realizations
\[
\mathfrak R_0,\mathfrak R_1,\ldots,\mathfrak R_T,
\]
let \(\mathfrak J_t:\mathfrak R_t\to\mathfrak R_{t+1}\) be admissible
correspondences. Define the path length
\[
\operatorname{Len}(\mathfrak J_{0:T})
=
\sum_{t=0}^{T-1}\operatorname{Cost}(\mathfrak J_t)
\]
and the induced directed path quasi-pseudometric
\[
\boxed{
d_{\mathrm{path}}^{\rightarrow}(\mathfrak R_0,\mathfrak R_T)
=
\inf_{\mathfrak J_{0:T}}
\operatorname{Len}(\mathfrak J_{0:T}).
}
\]
Concatenation gives the directed triangle inequality. Because replacement
costs and intervention divergences need not be symmetric, symmetry is not
assumed. When reverse correspondences exist, define the symmetrized
pseudometric
\[
\boxed{
d_{\mathrm{path}}^{\mathrm{sym}}(\mathfrak R,\mathfrak R')
=
\max\left\{
 d_{\mathrm{path}}^{\rightarrow}(\mathfrak R,\mathfrak R'),
 d_{\mathrm{path}}^{\rightarrow}(\mathfrak R',\mathfrak R)
\right\}.
}
\]
Distinct physical realizations can have zero symmetrized pseudodistance when
an exact correspondence exists in both directions. This is intentional:
identity is assigned to the invariant relational organization rather than to
the carrier tokens.

\subtopic{41.3 Identity failure modes}

A complete replacement audit must report at least
\[
\boxed{
\mathcal I_{\mathrm{rep}}
=
\left(
q_{\mathrm{str}},
\delta_{\mathrm{op}},
\delta_{\mathrm{causal}},
\delta_{\mathrm{rec}},
\mu^{\mathrm{keep}}
\right).
}
\]
The failure modes are different:
\[
\begin{array}{ll}
q_{\mathrm{str}}\text{ large}
&\text{the shared subspace changed},\\
\delta_{\mathrm{op}}\text{ large}
&\text{the same operations are no longer realized},\\
\delta_{\mathrm{causal}}\text{ large}
&\text{interventions no longer have matching consequences},\\
\delta_{\mathrm{rec}}\text{ large}
&\text{the new realization was not initialized from the old core},\\
\mu^{\mathrm{keep}}\text{ small}
&\text{content collapsed despite apparent coherence}.
\end{array}
\]
The theorem does not permit one successful coordinate to conceal failure of
another.
